\documentclass[11pt]{article}

\usepackage[margin=1in]{geometry}
\usepackage{amsmath,amssymb,amsthm}
\usepackage{graphicx}
\usepackage[hidelinks]{hyperref}

\newtheorem{theorem}{Theorem}
\newtheorem{lemma}{Lemma}

\newcommand{\D}{\mathbb D}
\newcommand{\C}{\mathbb C}
\newcommand{\B}{\mathcal B}
\newcommand{\ellp}{\ell^p}
\newcommand{\norm}[1]{\left\lVert #1\right\rVert}
\newcommand{\PsiNR}{\Psi}

\title{A Bohr--Bombieri Improvement of the Polynomial-Boundedness Constant}
\author{Automatic Math Discovery Run \texttt{fa\_banach\_001}}
\date{June 15, 2026}

\begin{document}
\maketitle

\section*{Status}

Status: \textbf{literature-implied answer (partial subcase)}.
The source paper asks whether the constant \(2\sqrt 3/3\) in its displayed
inequality (6.1) is optimal.  The answer is no: the classical
Bohr--Bombieri majorant formula lowers the constant to
\[
  B(1/2)=6-2\sqrt 6 < \frac{2\sqrt 3}{3}.
\]
This settles the optimality question for that displayed estimate, but it does
not answer the stronger question in the same paragraph asking for a norm-one
operator that is polynomially bounded with constant \(1\) and has
\(\PsiNR(T)=\infty\).

\section*{Source Question}

The source paper is H. Blazhko, D. Homza, F. L. Schwenninger,
J. de Vries, and M. Wojtylak, \emph{The algebraic numerical range as a
spectral set in Banach algebras}, arXiv:2410.10678.
In Example 6.3, after proving
\[
  \norm{f(T)}
  =\max\{\norm{f(E)}_2,\norm{f(\tfrac12 L)}_p\}
  \leq \frac{2\sqrt3}{3}\sup_{\D}|f|,
\]
the authors write:

\begin{center}
\includegraphics[width=\textwidth]{figures/open_problem_crop.png}
\end{center}

\section*{Supporting Literature}

Let \(f(z)=\sum_{n\ge0} a_n z^n\) be analytic on \(\D\) with
\(\norm{f}_{\infty,\D}\le 1\).  The Bohr--Bombieri function for the identity
operator is
\[
  B(r):=\sup_{\norm{f}_{\infty,\D}\le1}
        \sum_{n\ge0}|a_n|r^n .
\]
Bombieri and Ricci proved that, for \(1/3\le r\le 1/\sqrt2\),
\[
  B(r)=\frac{3-\sqrt{8(1-r^2)}}{r}.
\]
This formula is recalled explicitly, with citation to Bombieri--Ricci, in
Khasyanov's 2025 survey/article on majorant series estimates
\cite[Section 4]{Khasyanov2025}.  At \(r=1/2\) it gives
\[
  B(1/2)=6-2\sqrt6.
\]

\section*{Identification And Proof}

\begin{lemma}
\label{lem:shift-majorant}
Let \(L\) be the unilateral left shift on \(\ell^p\), \(1\le p\le\infty\).
For every polynomial \(f(z)=\sum_{k=0}^m a_kz^k\),
\[
  \norm{f(\tfrac12 L)}_p
  \leq (6-2\sqrt6)\sup_{|z|\le1}|f(z)|.
\]
\end{lemma}

\begin{proof}
Since \(\norm{L^k}_p\le1\) for every \(k\),
\[
  \norm{f(\tfrac12 L)}_p
  \leq \sum_{k=0}^m |a_k|2^{-k}.
\]
Applying the Bohr--Bombieri formula at \(r=1/2\), after normalizing by
\(\sup_{\D}|f|\), gives
\[
  \sum_{k=0}^m |a_k|2^{-k}
  \leq (6-2\sqrt6)\sup_{\D}|f|.
\]
\end{proof}

\begin{theorem}
\label{thm:not-optimal}
In Example 6.3 of arXiv:2410.10678, the constant \(2\sqrt3/3\) in inequality
(6.1) may be replaced by \(6-2\sqrt6\).  In particular, \(2\sqrt3/3\) is not
optimal for that estimate.
\end{theorem}

\begin{proof}
The Hilbert-space block \(E\) in the source example satisfies von Neumann's
inequality
\[
  \norm{f(E)}_2\le \sup_{\D}|f|.
\]
By Lemma~\ref{lem:shift-majorant},
\[
  \norm{f(\tfrac12L)}_p
  \leq (6-2\sqrt6)\sup_{\D}|f|.
\]
Therefore the direct-sum operator \(T=E\oplus \tfrac12L\), with the max
operator norm used in the source paper, satisfies
\[
  \norm{f(T)}
  = \max\{\norm{f(E)}_2,\norm{f(\tfrac12 L)}_p\}
  \leq (6-2\sqrt6)\sup_{\D}|f|,
\]
because \(6-2\sqrt6>1\).  Finally,
\[
  6-2\sqrt6 \approx 1.1010205
  < 1.1547005\approx \frac{2\sqrt3}{3}.
\]
Thus the displayed constant \(2\sqrt3/3\) is not optimal.
\end{proof}

\section*{Sharpness For The Shift Component}

For \(p=1\) or \(p=\infty\), the estimate in Lemma~\ref{lem:shift-majorant}
is sharp for the shift component.  Indeed, for \(p=1\) the matrix of
\(f(\tfrac12 L)\) has column sums tending to
\(\sum_{k=0}^m |a_k|2^{-k}\), and for \(p=\infty\) it has row sum equal to
the same expression.  Hence
\[
  \norm{f(\tfrac12 L)}_p=\sum_{k=0}^m |a_k|2^{-k}
  \qquad (p=1,\infty).
\]
The Bohr--Bombieri formula is sharp and its extremals are obtained by disk
automorphisms, approximable by polynomials because the relevant pole lies
outside \(\overline{\D}\).  Thus the constant \(6-2\sqrt6\) cannot be lowered
for the \(\ell^1\) or \(\ell^\infty\) shift block.

\section*{Scope And Limitations}

This packet answers only the optimality of the constant appearing in (6.1).
It does not produce a norm-one operator that is polynomially bounded with
constant \(1\) and has infinite numerical-range spectral constant.  The full
constant-\(1\) question remains open here.

\section*{Novelty And Search Notes}

Local duplicate check: the run registry, solution index, attempt index, and
proof-gap index were searched for arXiv:2410.10678, algebraic numerical range,
Bohr--Bombieri, polynomially bounded constant \(1\), and the constant
\(2\sqrt3/3\).  No existing packet for this target or improvement was found.

External/source check: arXiv:2410.10678 was inspected directly from its TeX
source and PDF.  The supporting Bohr--Bombieri formula was checked in a 2025
paper that explicitly cites Bombieri--Ricci's 1962 formula.  This is recorded
as a literature-implied answer rather than an original new proof of the
Bohr--Bombieri theorem.

\section*{Human Review Recommendation}

Reviewers should check only the import of the Bohr--Bombieri formula and the
standard computation of the \(\ell^1/\ell^\infty\) shift norm.  Once those are
accepted, the replacement of \(2\sqrt3/3\) by \(6-2\sqrt6\) is immediate.

\begin{thebibliography}{9}

\bibitem{Blazhko2024}
H. Blazhko, D. Homza, F. L. Schwenninger, J. de Vries, and M. Wojtylak,
\emph{The algebraic numerical range as a spectral set in Banach algebras},
arXiv:2410.10678.

\bibitem{BombieriRicci1962}
E. Bombieri and G. Ricci,
\emph{Sopra un teorema di H. Bohr e G. Ricci sulle funzioni maggioranti delle
serie di potenze},
Boll. Un. Mat. Ital. \textbf{17}(3), 276--282, 1962.

\bibitem{Khasyanov2025}
R. S. Khasyanov,
\emph{Area functional and majorant series estimates in the class of bounded
functions in the disk},
arXiv:2503.16313, Section 4.

\end{thebibliography}

\end{document}
